\documentclass[11pt,a4paper,fontset=fandol]{ctexart}
\usepackage[margin=22mm]{geometry}
\usepackage{xcolor,booktabs,tabularx,longtable,enumitem,listings,xurl,needspace,etoolbox}
\usepackage[colorlinks=true,linkcolor=blue!45!black,urlcolor=blue!55!black]{hyperref}
\usepackage{fancyhdr}
\pagestyle{fancy}\fancyhf{}\fancyhead[L]{STL Houjie Skills · 第一版}\fancyhead[R]{DeepSeek Harness}
\fancyfoot[C]{\thepage}\setlength{\headheight}{15pt}
\setlength{\parskip}{4pt}\setlength{\parindent}{2em}
\setlist{itemsep=3pt,topsep=4pt}
\lstset{basicstyle=\ttfamily\small,breaklines=true,columns=fullflexible,
keepspaces=true,showstringspaces=false,frame=single,rulecolor=\color{black!18},
backgroundcolor=\color{black!2},xleftmargin=2mm,xrightmargin=2mm,
aboveskip=8pt,belowskip=8pt,tabsize=2}
\emergencystretch=3em\sloppy
\setcounter{tocdepth}{1}
\BeforeBeginEnvironment{lstlisting}{\Needspace{5\baselineskip}}
\title{《STL源码剖析》技能包\\DeepSeek Harness 从安装到使用教程}
\author{WEveboy · Codex 协作整理}
\date{第一版 V1 / 1.0.0 · 2026年11月1日}
\begin{document}
\maketitle
\noindent\textbf{这份教程带你完成：}准备环境、安装仓颉、安装本技能、首次提问、运行十一项教学模型及归档审计、复现一个完整实验、检查结果并安全升级。

\noindent 项目：\url{https://github.com/WEveboy/STL_Houjie-skill}。
发行标签：\texttt{edition-v1.0.0}。本教程的发行下载链接在相应Release实际发布后生效；若资产缺失，请使用仓库源码，并以发布记录为准。
{\small\setlength{\parskip}{0pt}\tableofcontents}
\newpage
\section{先认识两种不同的工具}
仓颉 \texttt{cangjie-skill} 是作者端的蒸馏工具：读取长内容、提取候选、验证来源与执行条件，再编译技能。\texttt{stl-houjie} 是已经整理的STL成品技能：根据你的问题读能力卡、调用资源、解释源码与验证边界。\textbf{使用成品技能不以先装仓颉为硬前提。}教程从仓颉讲起，是为了使初次使用者也能理解来源与后续维护。

本项目研究侯捷书中早期SGI STL的设计，并提供原创的小规模教学模型和现代C++17契约实验。它不包含原书PDF或整章OCR，不是完整STL实现，也不是能证明任意容器正确的静态分析器。你提出具体问题时，仍应给出目标库、编译器、标准版本和最小代码。

\subsection{准备一个工作目录}
下面Windows PowerShell示例使用公开示例目录 \path{C:/SkillLab}，你可换到自己的工作盘。先打开终端，确认Python可用；本项目模型只使用Python标准库。仓颉脚本另需PyYAML，建议用独立虚拟环境隔离。
\begin{lstlisting}
python --version
python -m pip --version
New-Item -ItemType Directory -Force C:/SkillLab
Set-Location C:/SkillLab
\end{lstlisting}
需要Python 3.10或以上。若Windows仅提供\texttt{py}，可把本教程的\texttt{python}替换为\texttt{py -3}，先确认它指向同一个安装。macOS/Linux可使用\texttt{python3}，其目录写法按系统调整。不要把不同Python的pip混在一起。

建议保留三类目录：下载包、解压项目、自己的实验输出。这样升级不会覆盖输入和结果。安装目录只保存成品资源，不应当作日常实验仓库使用。
\section{从官方刷新包安装仓颉}
本次作者工具基线是仓颉2.5.0刷新包。官方说明指出，2026-09-13的刷新仍沿用版本号2.5.0，自动生成的旧tag源码ZIP不包含该刷新。应从上游README的刷新包入口或Release资产下载，检查\path{BUILD_INFO.json}；不要只下载一个\path{SKILL.md}。

\noindent 上游与版本页：
\begin{itemize}
\item \url{https://github.com/kangarooking/cangjie-skill}
\item \url{https://github.com/kangarooking/cangjie-skill/releases/tag/v2.5.0}
\end{itemize}
刷新通用包名为\path{cangjie-skill-2.5.0.zip}，旁边有同名\path{.sha256}校验文件。下载后先核对SHA256，再解压，找到完整\path{cangjie-skill/}目录；它应包含入口、\path{scripts/}、\path{methodology/}和\path{schemas/}等内容。不要误把外层解压目录当技能目录。

\begin{lstlisting}
Get-FileHash ./cangjie-skill-2.5.0.zip -Algorithm SHA256
Get-Content ./cangjie-skill-2.5.0.zip.sha256
\end{lstlisting}
逐字符比对两个哈希值，不比较文件名字段。SHA256说明下载字节与公布清单一致，不等于对发布者身份或内容安全作出独立认证。若不一致，停止使用该包，重新从确认的来源下载。

进入含\path{scripts/cangjie.py}的解压目录，创建作者端环境并做诊断：
\begin{lstlisting}
python -m venv .venv
./.venv/Scripts/python.exe -m pip install PyYAML
./.venv/Scripts/python.exe scripts/cangjie.py doctor
\end{lstlisting}
Linux/macOS把虚拟环境解释器路径改为\path{.venv/bin/python}。遇到明确列出的额外依赖缺失，按doctor提示在同一环境安装；可选依赖警告与必需依赖失败要分开读。不要为了“消除警告”无目标地安装整套开发环境。

成功检查应能说明工具版本及必需依赖状态。把输出保存在自己的维护记录中。接下来的安装位置由平台决定；不要把解压位置、宿主技能发现目录和本项目Release根目录混为一谈。
\section{把仓颉和STL成品接入DeepSeek Harness}
\subsection{先区分网页聊天与Harness}
普通DeepSeek网页对话不等于能读取本地能力卡并执行Python的Harness。这里的DeepSeek版是Harness插件，包含技能注册适配层及资源目录。使用真实模型时需要按官方流程自行配置模型连接；本项目不分发密钥，也没有替你运行付费模型会话。

官方项目：\url{https://github.com/deepseek-ai/deepseek-harness}。2026-11-01查阅的README将Harness标为开发预览，接口可能发生兼容性破坏。先安装适合该版本的Node.js和npm，再按官方入口启动：
\begin{lstlisting}
node --version
npm --version
npx @deepseek-ai/dsh web
\end{lstlisting}
终端保持运行，按其输出访问本地Web UI。不要把历史\texttt{rc.2}版本号当作今天的最新版。若希望使用本教程后续简短的\texttt{dsh}命令，可采用全局CLI方式：
\begin{lstlisting}
npm install -g @deepseek-ai/dsh
dsh --help
\end{lstlisting}
如果选择npx而未全局安装，后续相应调用也要使用npx形式；不要因找不到\texttt{dsh}就误判插件损坏。升级Harness后应重新检查插件API和注册契约，不能用一次旧版本测试替代新宿主验证。

\subsection{给同一个profile安装仓颉插件}
从仓颉v2.5.0 Release下载\path{dsh-cangjie-skill-2.5.0.tgz}及同名\path{.sha256}。按前面的SHA核验办法确认一致，再从本地TGZ安装。上游给出的web profile用法如下，启动也使用对应web入口：
\begin{lstlisting}
dsh plugin --profile web add C:/SkillLab/dsh-cangjie-skill-2.5.0.tgz
dsh web
\end{lstlisting}
profile不能一边装到web、一边在另一个隔离配置中寻找插件。仓颉脚本的Python作者环境仍需Python3.10以上与PyYAML；插件注册成功不证明这些外部命令也已安装。若你要重新蒸馏，请保留前节的完整刷新通用包作为作者工具，运行doctor，而非只依赖一个入口文本。

\subsection{安装STL第一版插件}
从本项目Release下载\path{stl-houjie-DeepSeek-1.0.0.zip}和\path{SHA256SUMS.txt}。ZIP便于取得安装说明和配套文件；原生插件资产是\path{dsh-stl-houjie-1.0.0.tgz}。两种形式选择需要的文件即可，不必把同一插件重复安装。
\begin{lstlisting}
Get-FileHash ./dsh-stl-houjie-1.0.0.tgz -Algorithm SHA256
Get-Content ./SHA256SUMS.txt
dsh plugin --profile web add C:/SkillLab/dsh-stl-houjie-1.0.0.tgz
dsh web
\end{lstlisting}
确认日志没有注册错误后新建任务，发送下一节提示词。适配层把技能与实际资源目录一起注册；如果宿主报告能看到技能却打不开相对卡片，应核对\texttt{resourceBase}目录契约、TGZ是否完整及宿主版本，而不是把卡片全文粘进系统提示掩盖问题。

\subsection{找到真实资源后再运行}
插件安装位置可能随Harness和profile变化，不在本教程中硬编码。请先让宿主显示已加载的\texttt{stl-houjie}来源及资源基目录，确认其中有\path{resources/stl_lab.py}。如果手动复现，下载本项目源码并在根目录运行后续命令更直接，不必修改宿主内部目录。

本项目对插件的模块导出、技能注册和资源定位做本地契约检查；真实模型会话没有执行，不能据此承诺账号配置、付费调用、宿主自动触发或具体模型结果已经实测成功。遇到这些问题，应保存宿主版本和脱敏日志，再按官方文档排查。
\section{第一次提问：同时给任务、输入与边界}
首次调用推荐显式写技能名，使“是否找得到入口”和“任务是否做对”更容易分开。可以先发送：
\begin{quote}
使用stl-houjie。按侯捷书中SGI模型解释vector为什么size与capacity不同；用size=3、capacity=4、在index=1插入2项推演状态。区分原书实现策略与现代标准保证，给出可复算步骤。
\end{quote}
如果宿主支持美元符号显式技能选择，也可写\texttt{\$stl-houjie}。关键是你看见宿主实际加载了正确入口和能力卡，而不是模型仅凭常识回答了同名主题。

一个有用的回答应明确初始有效范围、为何容量不足、新模型容量为6及新size为5，并说明这是所列教学策略，不是所有现代vector的增长倍率。它不能仅以“通常翻倍”替代输入条件。若你问真实源码，请继续提供库名、版本、代码和观测结果。

常用提问模板：
\begin{itemize}
\item “分析这段copy：标出源范围、目标范围、对象是否已存在、重叠方向，先判断前提再执行。”
\item “按给定比较器解释set/map如何确定键等价；不要把hash碰撞当相等。”
\item “画出deque的区号和区内偏移，B以元素计，解释负位移；不要把B当字节。”
\item “remove/unique以后哪些位置是有效结果？输出新尾，别把size仍未变说成错误。”
\end{itemize}
若问题缺少必要条件，技能应提出针对性补充，或将例子明确标为合成演练；不应猜出一个看似合理的源码版本。若要求运行代码，宿主需具备相应本地工具权限；纯聊天结果不等于已执行。
\section{从项目根目录运行十一项教学模型及归档审计}
以下均为本项目原创、有预算限制的整数教学模型。安装包\path{resources/usage.md}是输入合同的完整版本；此处给复制即可运行的起点。先找到发行项目根目录，再运行：
\begin{lstlisting}
python skills/stl-houjie/resources/stl_lab.py demo
\end{lstlisting}
每次CLI调用输出一个JSON。成功退出码0；输入、前提或预算错误为退出码2，且\texttt{status}为\texttt{input\_error}。\texttt{demo}一次展示8项操作；它不是8次模型对话，也不要求联网。

\subsection{run方式：一个文件对应一次操作}
把以下任意一段JSON保存为\path{input.json}，使用同一入口执行。用文本编辑器保存UTF-8即可；Windows UTF-8 BOM也受支持。不要把Markdown代码围栏写入JSON文件。
\begin{lstlisting}
python skills/stl-houjie/resources/stl_lab.py run input.json
\end{lstlisting}

\subsection{1. pool-class：小请求分档}
\begin{lstlisting}
{"op":"pool-class","bytes":30}
\end{lstlisting}
30字节按所列SGI参数上调到32字节，自由链表使用零下标3。128仍为小档，超过128走一级路径。模型不申请真实内存，也不模拟线程或任意对齐类型。

\subsection{2. vector-growth：大小和容量状态}
\begin{lstlisting}
{"op":"vector-growth","size":3,"capacity":4,
 "insertions":[{"index":1,"count":2}]}
\end{lstlisting}
新size为5；容量不足，教学重配置公式为旧size加max(旧size,插入数)，得到6。它不计算真实对象的复制/移动次数与异常保证，也不代表目标标准库的内部参数。

\subsection{3. deque-step：分段位移}
\begin{lstlisting}
{"op":"deque-step","buffer_elements":8,"offset":1,"delta":-3}
\end{lstlisting}
从区内1向前退3，到前一区的区内6。参数\texttt{buffer\_elements}以元素为单位；模型只给相对区位，不能证明该位置在某个真实deque里合法。

\subsection{4. heap-trace：维护最大堆}
\begin{lstlisting}
{"op":"heap-trace","values":[9,7,8,2],"action":"pop"}
\end{lstlisting}
输入必须已经是有效的零下标最大堆。pop把最大值9放到末尾，整体数组长度保持4，有效堆范围缩短为3。push操作另给\texttt{value}；不要把算法维护当成容器自动删除。

\subsection{5. logical-end：保留有效前缀}
\begin{lstlisting}
{"op":"logical-end","values":[1,1,2,1],"action":"unique"}
\end{lstlisting}
有效前缀是1、2、1，新尾为3，原size为4。unique只压缩相邻重复。若action改为remove，应另给\texttt{value}；不需要的字段不能保留凑数。

\subsection{6. bounds：有序范围的边界}
\begin{lstlisting}
{"op":"bounds","values":[1,2,2,4],"value":2}
\end{lstlisting}
非递减整数范围中，lower为1，upper为3，等价范围是[1,3)。模型不接受无序输入或任意用户比较器，不会静默排序后假装验证了原范围。

\subsection{7. hash-buckets：数学取模与重散列}
\begin{lstlisting}
{"op":"hash-buckets","values":[1,6,-1],"bucket_count":5,
 "new_bucket_count":7}
\end{lstlisting}
按非负数学余数分桶；重复值保留，改桶数后重新映射。它不仿真\texttt{std::hash}，不承诺SGI或现代库的实际遍历顺序，负数处理约定要和真实C++实现分开。

\subsection{8. copy-plan：同一数组中的合法复制}
\begin{lstlisting}
{"op":"copy-plan","values":[1,2,3,4,5],"first":0,"last":3,
 "result":5,"direction":"backward"}
\end{lstlisting}
此处result是目标尾点，目标[2,5)，结果为1、2、1、2、3。forward时result改为目标起点，不能只换direction而继续沿用同一个数字解释。非法重叠及目标对象不足会被拒绝，不执行未定义行为来观察输出。

输入预算为文件不超过64KiB，数组/操作数最多256项，模型储存和桶数最多4096，普通整数绝对值不超过1000000。程序拒绝重复JSON键、NaN/Infinity、未知字段以及布尔值冒充整数。预算或合同不满足时，应先修正任务，而非通过删除校验来运行。
\section{完整实验：合法右移、故障对照与保存结果}
目标是把已有数组前三项右移到索引[2,5)，保持仍需读取的源项不被提前覆盖。这个任务适合\texttt{copy-plan}及复制契约卡。它是自拟教学实验，不是原书某次运行的伪装复刻。

\subsection{建立输入与执行}
从项目根目录运行；PowerShell的单引号here-string使JSON不受变量替换影响。
\begin{lstlisting}
New-Item -ItemType Directory -Force experiments/copy-right
@'
{"op":"copy-plan","values":[1,2,3,4,5],"first":0,"last":3,
 "result":5,"direction":"backward"}
'@ | Set-Content -Encoding utf8 experiments/copy-right/input.json
python skills/stl-houjie/resources/stl_lab.py run experiments/copy-right/input.json > experiments/copy-right/output.json
\end{lstlisting}
读取\path{output.json}，核对最外层status为ok，结果对象中\texttt{values}为\texttt{[1,2,1,2,3]}，\texttt{destination}为\texttt{[2,5]}，\texttt{returned\_index}为2。结果之外还应读scope和contract，确认这是同一已有整数数组的有限模型。

\subsection{手算核对，不只检查“程序没报错”}
从尾向前赋值：先把索引2的3写到4，再把索引1的2写到3，最后把索引0的1写到2。每次读取的源仍未被覆盖，最终前两项保留1、2。返回索引2表示本次目标的起点；它不是删除后的容器新尾。

把方向改为forward并把result改成2，就是同一向右覆盖的危险方向：result落在源[0,3)内，模型应返回input\_error与退出码2。若保持result=5而改forward，则它成了从索引5开始写三项，失败原因是目标越界。两个失败条件不同，修复应对准其输入语义，不要混称为“copy都不能重叠”。

\subsection{归档最小可复现记录}
至少保留输入、完整输出、Python版本、资源文件哈希及你核对过的断言：
\begin{lstlisting}
python --version
Get-FileHash skills/stl-houjie/resources/stl_lab.py -Algorithm SHA256
Get-FileHash experiments/copy-right/input.json -Algorithm SHA256
Get-FileHash experiments/copy-right/output.json -Algorithm SHA256
\end{lstlisting}
哈希只帮助确认字节是否改变，不能代替上面的数学与合同核对。不要把凭据或私人绝对路径收进公开日志；与他人共享时，给项目版本、相对文件和必要环境就够了。
\section{再用真实C++与项目检查区分模型和实现}
Python模型有助于推演，却不证明某个C++类型的对象生存期或迭代器操作合法。源码项目包含原创C++17契约实验，需要本机可用的\texttt{g++}、\texttt{clang++}或\texttt{cl}。先让相应编译器在当前终端的PATH中可见；MSVC通常要在配置好的开发者终端使用。
\begin{lstlisting}
python scripts/run_cpp_checks.py
\end{lstlisting}
运行器记录实际编译器版本、编译与执行状态；默认结果保存在\path{build/cpp-checks/result.json}。若找不到源文件，可按\texttt{--help}指定\texttt{--source}和\texttt{--build-dir}。不建议双击CPP文件期待自动运行。

实验检查储存/构造分离、未初始化复制失败清理、traits与迭代器能力、vector重配置、list接合、deque访问、堆逻辑pop、map/哈希键、复制方向、逻辑新尾、二分边界、含重复项的nth\_element、引用修改、反向base与插入输出等。它保持断言启用，不固定现代库的容量倍率或桶数，不解引用失效迭代器。

本版作者工作区已完成的检查规模为Python 26组、C++ 16组、安装器6项。该数字只表示对应检查程序的范围，不是18项能力均经同一真实模型实测的证明。独立任务盲测、路由评测和来源审计以仓库随版报告为准；本教程不预填尚未得到的结果，也不声称比不使用Skill更准确。

从源码根目录可重新运行Python和安装器测试：
\begin{lstlisting}
python -m unittest discover -s skills/stl-houjie/resources -p test_stl_lab.py -v
python -m unittest discover -s tests -v
\end{lstlisting}
若只有精简安装包而没有\path{tests/}或\path{scripts/}，先下载源码项目，不要把“资源没有随这个包提供”误判为算法失败。运行成功后保存日志及版本，失败则先读第一处错误和当前环境，不用删除断言掩盖问题。
\section{怎样读这本书的历史边界}
先读\path{source/CHAPTER_READING.md}建立八章关系，再按任务查\path{source/DIGEST.md}和\path{skills/stl-houjie/references/capabilities/}。一次具体操作可按“储存哪里来—对象何时存在—迭代器能做什么—算法要求什么—策略与配接如何改变行为”追踪。

\begin{itemize}
\item SGI的小块128字节阈值、8字节分档、vector增长公式、deque区容量及哈希质数表都带实现限定，不能写成现代所有标准库必须采用的常量。
\item 已分配储存不等于已有目标对象；reserve后size仍为0，不能据此向begin普通copy若干元素。
\item remove/unique给逻辑新尾而不自动缩小容器；pop\_heap不删除末项；unique只处理相邻重复。
\item deque中的随机访问不保证整个序列是连续地址；反向迭代器的base不等于它正在解引用的位置。
\item 原书deque部分单位与算式、nth\_element描述、for\_each简化分类有局部冲突，按来源审计分别处理；不要为“忠于原文”继续执行错误条件。
\item 书中部分异常及线程路径被省略，不能从节录推导完整强保证或线程安全。完整红黑删除修复不在本版的原书插入讲解保证内。
\end{itemize}
若你研究MyTinySTL或另一现代库，先提供该仓库对应源码与提交；本技能能帮助建立分析框架，不能因为看过侯捷书就宣布另一个实现正确。现代接口及移除年代等事实需查当前官方标准或目标库资料，并与原书证据分开。
\section{常见故障：从最小可确认的状态开始}
\begin{description}[style=nextline]
\item[找不到python或pip] 先重开终端并确认安装与PATH；用\texttt{python -m pip}避免pip属于另一解释器。Python模型无第三方依赖，只有仓颉作者端另需PyYAML。
\item[技能没有出现] 检查是否整目录安装、入口是否多嵌一层、同名入口是否重复。Codex检查所选发现位置；Harness检查profile及注册日志。宿主未实际加载不能用模型口头声称替代。
\item[stl\_lab.py路径不存在] 先定位完整成品根目录。源码在\path{skills/stl-houjie/resources/}，安装版在安装器输出目标下，Harness以注册的resourceBase为准。
\item[JSON返回input\_error] 先读message。检查字段名、整数与布尔值、范围、排序及堆序前提。UTF-8 BOM受支持，代码围栏、单引号JSON和末尾多逗号不是有效输入。
\item[得到和书中不同的capacity或bucket数] 先核所用实现。教学模型、历史SGI输出和你本机标准库的参数不是同一个承诺；不要为了匹配书图修改测试目标。
\item[C++实验找不到编译器] 在配置好的开发者终端运行，检查编译器版本。返回unavailable不能写成16组已在你的机器通过；仅靠Python通过也不能补上这项执行。
\item[插件加载后打不开资源] 检查TGZ是否完整、注册时目录资源契约及宿主版本。不要只复制\path{index.mjs}；它依赖同包的技能及资源目录。
\item[模型没有实际运行工具] 请求它报告命令、退出码和产物位置；若宿主无本地工具，就自己按教程命令运行。模型回答的预期结果与实际输出应分开保存。
\end{description}
\section{升级、备份与宿主兼容}
升级前保存当前Harness版本、profile、插件TGZ及其SHA和本项目版本，并备份自己的实验文件。先在测试profile或单独环境验证新宿主与插件；具体移除/替换语法以该版本\texttt{dsh --help}和官方插件说明为准，本教程不把开发预览时期的未知升级接口写成固定保证。

新TGZ仍先校验后安装，保持profile一致，并查看是否有同名重复注册。用首个提示词、资源读取和一次demo检查替换后的实际状态。若新版本不兼容，保留错误日志，按宿主支持的管理方式恢复旧TGZ；不修改密钥或删除整个profile来试错。

本项目只带自己的插件适配层和原创资源，上游Harness、仓颉及模型依赖由用户按其许可和版本自行安装。真实付费会话未执行；后续宿主升级后需要重新做注册契约和端到端调用验证，不能继续引用旧测试结论当当前保证。
\section{继续精读、重新蒸馏与维护}
若只想使用现有能力，到此即可。若想加入新的源码版本或补全未覆盖机制，先保留旧版，使用仓颉整理新增证据，不手改编译后入口冒充源数据更新。你可向作者端宿主发送：
\begin{quote}
用cangjie-skill继续维护本项目。先读source中的整书骨架、覆盖审计与流水线状态，再核对新材料。保留稳定能力ID；只修改有来源或实际失败支持的卡片。分别记录来源核查、代表任务输出和边界测试，重新编译后再打包。
\end{quote}
本项目的作者源与可编译Bundle同在\path{source/}，其中含verified.yaml、cards、book与resources；不要把上游旧示例的侧车路径套到本发行目录。重新蒸馏需要合法可读的原始材料，公开仓库并不附带整本书。不能只凭模型记忆补出原书没有写的步骤或把现代补充归给作者。

维护循环为：更新关键任务与候选、检查来源充分性/可执行性/任务增益、回填RIA能力卡、补输出与边界评测、编译single成品、验证安装与两平台资源一致性、记录版本及发布。阅读稿负责解释全书，能力卡负责完成任务，模型脚本负责有限可复算部分，三者不互相替代。

发布前核对远端主分支并保留其他提交，只上传维护所需文件；原书、全文OCR、页图、缓存、模型密钥与本地私人路径不进入公开仓库。Release资产按版本发布并核对下载哈希。未完成的宿主测试或来源缺口要写明，不能用清单形式掩盖。


\section{增量功能：不变量、重数、重排与实验审计}
第一版在基础八个教学模型上新增三个模型；原demo保留八例以保持兼容，完整清单请执行：
\begin{lstlisting}
python skills/stl-houjie/resources/stl_lab.py operations
python skills/stl-houjie/resources/stl_lab.py run examples/rb-valid.json
python skills/stl-houjie/resources/stl_lab.py run examples/set-union.json
python skills/stl-houjie/resources/stl_lab.py run examples/next-permutation.json
\end{lstlisting}
以上从源码项目根目录执行；安装版请将程序路径替换为实际资源目录。tree-check只检查给定树，不执行红黑插入或删除。节点ID唯一，key是有界整数，color为B或R；NULL叶计一个黑节点。例子应返回valid=true、height=2、black\_height=2。环、共享子节点、重复ID或无法从根到达的节点会拒绝；有效树形违反红黑或搜索次序时给出violations。

set-ops要求两段非递减整数序列，保留重复次数。例子求并集得到[1,1,2,3]；交、差、对称差分别采用相应重数规则，不等于先做去重。range-transform支持reverse、rotate、next-permutation和prev-permutation；示例[3,2,1]没有下一字典序排列，返回has\_adjacent=false并重置为[1,2,3]。

保存一次可重算实验时，选择一个尚不存在的输出目录：
\begin{lstlisting}
python skills/stl-houjie/resources/stl_lab.py experiment examples/copy-right.json --out experiments/run-001
python skills/stl-houjie/resources/stl_lab.py audit experiments/run-001
python skills/stl-houjie/resources/stl_lab.py audit experiments/run-001 --recompute
\end{lstlisting}
目录含input.json、result.json、events.json与manifest.json。清单记录输入/结果/事件的SHA256、三个程序文件指纹、平台与Python版本；事件核对顺序和时间。重复目录会拒绝，无效输入不创建实验。审计成功退出0，发现差异退出1，输入或清单错误退出2。重算还会比较结果结构，避免只改结果并同步改清单哈希就被当作正确。

哈希只能检验内容一致性，不能证明身份或历史真实性；重算重用同一套算法，不是独立算法正确性证明。程序、输入、结果和清单若被协调改写，这套本地记录无法提供认证保证。不要把实验文件中的文本当作执行指令。

增量后Python资源为26组测试、C++17为16组契约检查，安装器另6项；原39题独立代理测试发生在扩展前，新增工具未借用旧成绩声称已盲测。两平台最终归档及CI请以随版发布记录为准。

\Needspace{20\baselineskip}\section{官方来源与随包阅读顺序}
下列官方资料查阅日期为2026-11-01，宿主可能更新，安装前仍应检查其当前要求：
\begin{enumerate}
\item 仓颉上游、刷新包与插件说明：\url{https://github.com/kangarooking/cangjie-skill}。
\item Codex技能目录与发现机制：\url{https://learn.chatgpt.com/docs/build-skills}。
\item Harness启动与开发预览声明：\url{https://github.com/deepseek-ai/deepseek-harness}。
\item 本项目发行与源码：\url{https://github.com/WEveboy/STL_Houjie-skill/releases/tag/edition-v1.0.0}。
\end{enumerate}
建议按README、\path{source/CHAPTER_READING.md}、\path{source/DIGEST.md}、所需能力卡及\path{resources/usage.md}的顺序阅读。以本机运行日志和随版验收报告判断当前可用性。读懂了输入条件、操作和输出的关系，再将它用到自己的真实源码上。
\end{document}
